\documentclass[10pt]{article}
\lineskip=0pt

\usepackage[skins,breakable,raster]{tcolorbox}
\tcbsetforeverylayer{
  skin=enhanced jigsaw,
  fonttitle=\bfseries\small,
  parbox=false,
  colback=white,
  subtitle style={colback=gray!15,fontupper=\scriptsize\bfseries},
  left=1pt,right=1pt,top=1pt,bottom=1pt,
  boxrule=0.35pt,
  nobeforeafter,
  % multicol compatibility
  breakable, height fixed for=first,
  pad at break*=0pt, pad before break*=0pt, pad after break=0pt,
  break at=-\baselineskip/0pt
}

\usepackage{parskip}
\setlength\parskip{0pt}
\usepackage{multicol}
\usepackage{geometry}
\newgeometry{noheadfoot=true,tmargin=5pt,bmargin=5pt,hmargin=5pt}
\geometry{paperheight=297mm,paperwidth=420mm}
\usepackage{xcolor}
\usepackage{soul}
\usepackage[fleqn]{amsmath}
\usepackage{amssymb,mathtools}
\usepackage[inline]{enumitem}
\usepackage{tabularx}
\usepackage{tabularray}
\usepackage{graphicx}
\pagestyle{empty}

\setlength{\parindent}{0pt}
\setlength{\columnsep}{2pt}
\setlist{leftmargin=*,labelsep=2pt,itemsep=0pt,topsep=0pt,parsep=0pt}
\setlength{\abovedisplayskip}{1pt}
\setlength{\belowdisplayskip}{1pt}
\setlength{\abovedisplayshortskip}{0pt}
\setlength{\belowdisplayshortskip}{0pt}
\newcommand{\warn}[1]{\textbf{\color{red!65!black}#1}}
\newcommand{\key}[1]{\textbf{#1}}
\newcommand{\ind}[1]{\mathbf 1_{\{#1\}}}
\newcommand{\boxstart}[1]{\begin{tcolorbox}[title={#1},sharp corners]}
\newcommand{\boxend}{\end{tcolorbox}}

\begin{document}
\small

% PAGE 1: probability, conditioning, random variables, indicators
\begin{multicols*}{3}
	\vspace*{-1em}

	\boxstart{L1 — How to attack a probability question}
	\begin{enumerate}
		\item Define the experiment and sample space so outcomes are equally likely (if using counting).
		\item Translate the wording into events; write the complement for ``at least one'' or ``none''.
		\item Decide whether stages are independent, conditional, with/without replacement, or disjoint.
		\item Count by stages, or draw a probability tree and multiply along paths.
		\item Add disjoint cases; check that all cases form a partition and probabilities sum to 1.
	\end{enumerate}
	\textbf{Counting patterns:} choose the pattern first, then choose values, then arrange positions. For repeated values, divide by the factorials of identical groups. If a choice is made as a set, do not treat it as ordered.

	\smallskip
	\textbf{Sanity checks:} probabilities lie in $[0,1]$; a partition counts every outcome once; a conditional probability is computed inside the stated information; ``at least one'' is often easiest by complement.
	\boxend

	\boxstart{L1 — Collision / birthday workflow}
	For $k$ independent assignments into $m$ equally likely slots:
	\[
		P(\text{no collision})=\prod_{i=0}^{k-1}\frac{m-i}{m},
		\qquad P(\ge1\text{ collision})=1-P(\text{none}).
	\]
	Use the exact product for small values; approximate by $1-e^{-k(k-1)/(2m)}$ when appropriate. A collision among any pair needs about $\sqrt m$ trials; matching one fixed target needs order $m$ trials. Do not confuse these questions.
	\boxend

	\boxstart{L1 — Lecture pattern: counting a partition}
	For a five-digit string over $\{0,\ldots,9\}$, classify by repeat pattern before counting. For example,
	\[
		\begin{array}{c|c}
			\text{pattern}         & \text{counting idea}                     \\ \hline
			\text{all distinct}    & P(10,5)                                  \\
			\text{one pair}        & \binom52\,10\,P(9,3)                     \\
			\text{two pairs}       & \binom52\binom32\binom{10}{2}\,P(8,1)/2! \\
			\text{three of a kind} & \binom53\,10\,\binom92\,2!               \\
		\end{array}
	\]
	The exact remaining patterns follow the same principle. Check that the mutually exclusive pattern counts sum to $10^5$. The division by $2!$ prevents ordering identical groups twice.
	\boxend

	\boxstart{L1 — Uniform shuffle: Fisher--Yates}
	For $i=n-1,n-2,\ldots,1$, choose $j$ uniformly from $\{0,\ldots,i\}$ and swap positions $i$ and $j$. The choice counts are $n(n-1)\cdots2=n!$, and each choice sequence maps bijectively to one permutation, so every output has probability $1/n!$.

	\warn{Pitfall:} choosing $j$ uniformly from the full range $\{0,\ldots,n-1\}$ at every step gives $n^n$ equally likely choice sequences, which generally do not map evenly to the $n!$ permutations.
	\boxend

	\boxstart{L1 — Union bound vs inclusion--exclusion}
	For events $A_1,\ldots,A_n$,
	\[
		P\!\left(\bigcup_iA_i\right)\le\sum_iP(A_i).
	\]
	The general inclusion--exclusion pattern is
	\[
		P\!\left(\bigcup_{i=1}^nA_i\right)
		=\sum_{k=1}^n(-1)^{k+1}
		\sum_{i_1<\cdots<i_k}
		P(A_{i_1}\cap\cdots\cap A_{i_k}),
	\]
	that is, add singles, subtract pairs, add triples, and continue alternating. The bound is useful when the right side is small and intersections are rare; it is loose when events overlap heavily.

	For a correct shuffle, let $A_i$ mean item $i$ stays fixed. For any specified set of $k$ items, the remaining $n-k$ items can be permuted freely, so
	\[
		P(A_{i_1}\cap\cdots\cap A_{i_k})
		=\frac{(n-k)!}{n!}.
	\]
	There are $\binom nk$ choices of the $k$ fixed items. Hence the complete $k$-event contribution is
	\[
		\binom nk\frac{(n-k)!}{n!}=\frac1{k!}.
	\]
	Therefore the exact shuffle-test failure probability is
	\[
		P(\text{some fixed point})=1-\frac1{2!}+\frac1{3!}-\cdots+\frac{(-1)^{n+1}}{n!}\approx1-e^{-1}.
	\]
	The union bound alone only gives $P(\cup_iA_i)\le n(1/n)=1$.
	\boxend

	\boxstart{L2 — Conditioning and Bayes: path method}
	\begin{itemize}
		\item Draw one path for each mutually exclusive cause and evidence combination.
		\item Assign each path its prior/previous probability times the next conditional probability.
		\item Add all paths producing the observed evidence.
		\item Posterior for a cause $=$ its evidence-path weight divided by the total evidence weight.
	\end{itemize}
	\[
		\text{posterior odds}=\text{prior odds}\times\text{likelihood ratio}.
	\]
	\warn{Base rates matter: sensitivity is not $P(\text{cause}\mid\text{positive})$.}

	\textbf{Conditional independence:} $A\perp B\mid D$ means factorisation only after fixing $D$. It does not imply $A\perp B$ overall; different causes can induce marginal dependence.
	\boxend

	\boxstart{L2 — Independence traps}
	\begin{itemize}
		\item Independent means learning one event does not change the other’s probability.
		\item Disjoint positive-probability events are dependent, not independent.
		\item For random variables, check both factorisation of the joint function and product support. A factorised expression on a restricted support is not enough.
		\item \textbf{Pairwise independence:} every pair factors.
		\item \textbf{Mutual independence:} every finite subcollection factors; mutual implies pairwise, but not conversely.
		\item \textbf{Conditional independence:} $X\perp Y\mid Z$ means the factorisation holds after fixing $Z=z$; it may fail marginally because mixing over $Z$ induces dependence.
	\end{itemize}
	\textbf{XOR warning:} fair independent $X,K$ and $C=X\oplus K$ give pairwise independence but not mutual independence. \textbf{Naive Bayes:} within each fixed class, multiply feature likelihoods; multiply by the class prior; normalise over classes.
	\boxend

	\boxstart{L2 — Lecture pattern: fair extraction}
	For iid biased bits with $P(1)=p$, pair consecutive bits:
	\[
		01\mapsto0,\qquad10\mapsto1,\qquad00,11\mapsto\text{discard}.
	\]
	The accepted outputs are fair because $P(01)=P(10)=p(1-p)$. This fails if the two bits are not identically distributed or are dependent. Equal pairs can be recycled recursively by encoding $00$ and $11$ as a new biased stream.
	\boxend

	\boxstart{L2 — Uniform fixed-size streaming sample}
	To choose exactly $n$ of $N$ records uniformly while reading them once, keep state $(t,m)$ = (records remaining, selections still needed):
	\begin{enumerate}
		\item keep the next record with probability $m/t$;
		\item if kept, set $m\leftarrow m-1$; always set $t\leftarrow t-1$;
		\item stop when $m=0$; if $m=t$, keep every remaining record.
	\end{enumerate}
	At every state, $P(\text{keep})=m/t$ and $P(\text{skip})=(t-m)/t$. Multiplying path probabilities shows every $n$-subset has probability $1/\binom Nn$.
	\boxend

	\boxstart{L3 — Random-variable interpretation}
	A random variable maps outcomes to numbers. Always identify its support and what is being counted (e.g. failures before success versus trials including success).

	\textbf{PMF:} masses must be nonnegative and sum to 1. \quad\\
	\textbf{PDF:} nonnegative, total area 1; density height is not a probability.\\
	\textbf{CDF:} nondecreasing, right-continuous, limits $0$ and $1$. For a discrete variable, jumps are point masses; for a continuous variable, interval endpoints do not matter.

	\textbf{Unknown constant workflow:}
	\begin{enumerate}[label=\arabic*.,leftmargin=*,itemsep=0pt]
		\item Identify the support, i.e. determine all values that the random variable can take, like a domain.
		\item Normalise: use $\sum_x f_X(x)=1$ for a PMF, or $\int f_X(x)\,dx=1$ for a PDF.
		\item Substitute the constant back into the PMF/PDF.
		\item Sum or integrate the requested event; check nonnegativity and total mass/area $1$.
	\end{enumerate}
	\textbf{Example:} if $f_X(x)=cx$ for $0<x<1$,
	\[
		\int_0^1cx\,dx=1\ \Rightarrow\ c=2,
		\qquad P(X\le1/2)=\int_0^{1/2}2x\,dx=\frac14.
	\]
	\boxend

	\boxstart{L3 — Indicators and symmetry}
	For an event $A$, $I_A=\ind{A}$ is 1 if $A$ occurs and 0 otherwise. For a count $N=\sum_i I_i$,
	\[
		E(N)=\sum_iP(A_i),\qquad
		V(N)=\sum_iV(I_i)+2\sum_{i<j}\operatorname{Cov}(I_i,I_j).
	\]
	Use indicators when the full PMF is awkward. For symmetric positions, compute one position’s probability and multiply. For dependent indicators, determine joint probabilities for pairs rather than assuming independence.

	\textbf{Examples:} record count in a random permutation; descents/runs; number of occupied boxes; number of successful transmissions.
	\boxend

	\boxstart{L3 — Expectation tricks}
	\begin{itemize}
		\item Use $E[g(X)]$ directly when the question asks for a transformed quantity; do not first find the distribution of $g(X)$ unless needed.
		\item Linearity works without independence.
		\item Tail-sum method: for nonnegative integer $X$, $E(X)=\sum_{k\ge1}P(X\ge k)$; for nonnegative continuous $X$, $E(X)=\int_0^\infty P(X>t)\,dt$.
		\item Variance needs second moments or covariances; unlike expectation, it is not generally linear.
	\end{itemize}
	\textbf{Inverse-CDF generation:} if $U\sim U(0,1)$ and $F$ is continuous and increasing, set $X=F^{-1}(U)$. Then
	\[
		P(X\le x)=P(F^{-1}(U)\le x)=P(U\le F(x))=F(x),
	\]
	so $X$ has the target CDF (and PDF $f_X=F'$ where differentiable). For a non-strict CDF use the generalised inverse $F^{-1}(u)=\inf\{x:F(x)\ge u\}$.
	\boxend

	\boxstart{L3 — Finite binary counter: tail-sum method}
	Choose an $m$-bit counter state uniformly, add 1 with wraparound, and let $X$ be the number of flipped bits. At least $k$ bits flip exactly when the last $k-1$ bits are all 1:
	\[
		P(X\ge k)=2^{-(k-1)},\qquad k=1,\ldots,m.
	\]
	Using the integer tail-sum formula,
	\[
		E(X)=\sum_{k=1}^mP(X\ge k)
		=1+\frac12+\cdots+\frac1{2^{m-1}}
		=2-2^{-m}.
	\]
	This is a standard pattern: find simple tail events instead of constructing the full PMF.
	\boxend

	\boxstart{L3 — Inverse-CDF example: Pareto}
	For $f_X(x)=2/x^3$ on $x\ge1$,
	\[
		F_X(x)=1-\frac1{x^2},\qquad U=F_X(X)
		\quad\Longrightarrow\quad
		X=F_X^{-1}(U)=\frac1{\sqrt{1-U}}.
	\]
	Derive the CDF of the generated value to verify the simulation; do not only compare the shape of a histogram. Here $E(X)=2$ while the median is $\sqrt2$, illustrating the effect of a heavy upper tail.
	\boxend

	\boxstart{L3 — Complete collection / coupon collector}
	For $d$ equally likely independent designs, let $W_j$ be the number of boxes needed to move from $j-1$ collected designs to $j$. The next box is new with probability
	\[
		p_j=\frac{d-j+1}{d},\qquad W_j\sim\operatorname{Geom}(p_j),
		\qquad E(W_j)=\frac1{p_j}=\frac d{d-j+1}.
	\]
	Since $T=\sum_{j=1}^dW_j$, linearity gives
	\[
		E(T)=d\sum_{r=1}^d\frac1r=dH_d.
	\]
	This assumes equal design probabilities; a rare special design can dominate the waiting time and invalidate $dH_d$.
	\boxend

	\boxstart{L3 — Record indicators}
	For a uniformly random permutation, $R_i=\ind{\text{position $i$ is larger than all earlier positions}}$. The largest of the first $i$ values is equally likely to be in any position, so $P(R_i=1)=1/i$. Therefore, for $M=\sum_{i=1}^nR_i$,
	\[
		E(M)=H_n=\sum_{i=1}^n\frac1i.
	\]
	The lecture also shows $R_i,R_j$ are pairwise independent for $i<j$, giving $V(M)=\sum_i(1/i)(1-1/i)$.
	\boxend

	\boxstart{L1–L3 — Common wording translations}
	\begin{tabularx}{\linewidth}{@{}lX@{}}
		at least one               & use the complement of none                            \\
		exactly $k$                & identify the $k$ successes and all remaining failures \\
		first success on trial $k$ & first $k-1$ fail, then success                        \\
		$k$th success on trial $n$ & first $n-1$ contain $k-1$ successes, then success     \\
		given / among              & restrict to the conditioning event and renormalise    \\
		rate $r$ for time $t$      & expected count is $rt$                                \\
	\end{tabularx}
	\boxend

\end{multicols*}

\newpage

% PAGE 2: joint distributions, dependence, models, approximations, simulation
\begin{multicols*}{3}
	\vspace*{-1em}

	\boxstart{L4 — Joint-distribution workflow}
	\begin{enumerate}
		\item Draw or describe the support before integrating or summing.
		\item Normalise the joint function over the support.
		\item Marginalise by summing/integrating out the unwanted coordinate.
		\item Obtain a conditional distribution by dividing the joint function by the relevant marginal and renormalising on the remaining support.
		\item Compute events as sums/areas over the correct region, not rectangular bounds unless the support is rectangular.
	\end{enumerate}
	\textbf{Support warning:} a formula can factor but variables still be dependent if the support couples $x$ and $y$ (triangle, diagonal, wedge, fixed total, etc.).
	\boxend

	\boxstart{L4 — Dependence diagnostics}
	\begin{itemize}
		\item Independence: joint $=$ product of marginals everywhere on a product support.
		\item Independence $\Rightarrow$ zero covariance, but zero covariance $\not\Rightarrow$ independence.
		\item Covariance measures linear association, not arbitrary dependence or causation.
		\item Fixed-size sampling creates negative dependence: if one item is selected, another is slightly less likely to be selected.
	\end{itemize}
	\textbf{Counterexample pattern:} let $W$ be symmetric around 0 and set $Y=W^2$. Then covariance can be zero even though $Y$ is completely determined by $W$.
	\boxend

	\boxstart{L4 — Joint-table exam workflow}
	For a discrete joint table: sum columns/rows to get marginals, then divide the relevant row/column by its marginal to get a conditional distribution. For example,
	\[
		P(Y=y\mid X=x)=\frac{P(X=x,Y=y)}{P(X=x)},
		\qquad
		E(Y\mid X=x)=\sum_y yP(Y=y\mid X=x).
	\]
	To test independence, compare every cell with the product of its marginals. If independent, conditioning does not change the other variable’s distribution and $E(Y\mid X=x)=E(Y)$.
	\boxend

	\boxstart{Additional tutorial -- Q4: Poisson split given the total}
	If $X_1\sim\operatorname{Poisson}(\lambda)$ and $X_2\sim\operatorname{Poisson}(\mu)$ are independent, set $S=X_1+X_2$. Given $S=n$, $X_1$ can only be $0,1,\ldots,n$. For $x$ in this range,
	\[
		\begin{aligned}
			P(X_1=x\mid S=n)
			 & =\frac{P(X_1=x,\,X_2=n-x)}{P(S=n)}                        \\
			 & =\frac{e^{-\lambda}\lambda^x/x!\;e^{-\mu}\mu^{n-x}/(n-x)!}
			    {e^{-(\lambda+\mu)}(\lambda+\mu)^n/n!} \\
			 & =\frac{n!}{x!(n-x)!}\frac{\lambda^x\mu^{n-x}}
			                       {(\lambda+\mu)^n}              \\
			 & =\binom nx\left(\frac{\lambda}{\lambda+\mu}\right)^x
			\left(\frac{\mu}{\lambda+\mu}\right)^{n-x}.
		\end{aligned}
	\]
	So $X_1\mid(S=n)\sim\operatorname{Bin}\!\left(n,\frac{\lambda}{\lambda+\mu}\right)$ and
	\[
		E(X_1\mid S=n)=n\frac{\lambda}{\lambda+\mu}.
	\]
	\textbf{Tutorial 5:} $\lambda=2,\mu=3,n=10$ gives $\operatorname{Bin}(10,2/5)$ and answer $10(2/5)=4$.
	\boxend

	\boxstart{L4 — Runs in a random permutation}
	For continuous distinct values $A_1,\ldots,A_n$, let $D_i=\ind{A_i>A_{i+1}}$. A run starts at position 1 and after every descent:
	\[
		R=1+\sum_{i=1}^{n-1}D_i.
	\]
	Each $D_i$ has probability $1/2$. Non-adjacent descents are independent, while
	\[
		P(D_i=D_{i+1}=1)=\frac16,qquad
		\operatorname{Cov}(D_i,D_{i+1})=\frac16-\frac14=-\frac1{12}.
	\]
	Therefore, for $n\ge2$,
	\[
		E(R)=\frac{n+1}{2},\qquad V(R)=\frac{2n-1}{12}.
	\]
	\boxend

	\boxstart{L2/L4 — Fixed-size sampling dependence}
	If exactly $n$ of $N$ records are selected uniformly and $I_i$ indicates that record $i$ is included,
	\[
		P(I_i=1)=\frac nN,\qquad
		P(I_i=1,I_j=1)=\frac{n(n-1)}{N(N-1)}\quad(i\ne j).
	\]
	Thus selected indicators are negatively correlated: selecting one consumes one of the fixed slots. This is a standard example where marginal probabilities alone hide dependence.
	\boxend

	\boxstart{L5 — Reservoir sampling}
	To maintain a uniform sample of $k$ records after $i>k$ arrivals: draw $J_i$ uniformly from $\{1,\ldots,i\}$; if $J_i\le k$, replace slot $J_i$ by record $i$, otherwise do nothing. The new record enters with probability $k/i$.

	For replacements from indices $m+1$ through $n$, $R=\sum I_i$ has independent decisions with
	\[
		E(R)=k\sum_{i=m+1}^n\frac1i,qquad
		V(R)=E(R)-k^2\sum_{i=m+1}^n\frac1{i^2}.
	\]
	The decisions are independent; the reservoir contents are not independent observations.
	\boxend

	\boxstart{L5 — Choosing a discrete model}
	\begin{tabularx}{\linewidth}{@{}lX@{}}
		Bernoulli         & one success/failure trial                                         \\
		Binomial          & number of successes in fixed $n$ iid trials                       \\
		Geometric         & trials through the first success; memoryless waiting time         \\
		Negative binomial & trials through the $k$th success                                  \\
		Poisson           & event count in a fixed interval/region with constant rate         \\
		Poisson-binomial  & sum of independent Bernoulli variables with unequal probabilities \\
	\end{tabularx}
	\smallskip
	\textbf{Checklist:} fixed or random number of trials? identical success probability? independence? count or waiting time? If a rate is stated, rescale it to the observation window.
	\boxend

	\boxstart{L5 — Poisson-binomial distribution}
	If $I_1,\ldots,I_n$ are independent Bernoulli variables with possibly unequal success probabilities $p_i$, then
	\[
		R=\sum_{i=1}^n I_i,\qquad
		\lambda=E(R)=\sum_i p_i,
		\qquad
		V(R)=\sum_i p_i(1-p_i)=\lambda-\sum_i p_i^2.
	\]
	This is Poisson-binomial, not ordinary binomial unless all $p_i$ are equal. For $Z\sim\operatorname{Poisson}(\lambda)$ and any set $A$ of nonnegative integers,
	\[
		\left|P(R\in A)-P(Z\in A)\right|
		\le \frac{\sum_i p_i^2}{\lambda}
		\le \max_i p_i.
	\]
	Thus the Poisson approximation is strongest when every individual success probability is small; the largest $p_i$ controls the simple error bound.
	\boxend

	\boxstart{L5 — Nonlinear expected-value template}
	If a quantity depends nonlinearly on a random count, substitute moments rather than plugging in the mean:
	\[
		E(a t-cX^2)=a t-cE(X^2)
		=a t-c\{V(X)+[E(X)]^2\}.
	\]
	For $X\sim\operatorname{Poisson}(rt)$, this becomes
	\[
		E(a t-cX^2)=a t-c\{rt+(rt)^2\}.
	\]
	Then optimise the resulting deterministic function in $t$ and check it is a maximum. This is the method behind Tutorial 5 Question 3.
	\boxend

	\boxstart{L6 — Choosing a continuous model}
	\begin{tabularx}{\linewidth}{@{}lX@{}}
		Uniform     & every point in an interval equally likely; use length ratios         \\
		Exponential & nonnegative waiting time with constant hazard / Poisson interarrival \\
		Normal      & symmetric measurement/noise model; standardise before using $\Phi$   \\
	\end{tabularx}
	\smallskip
	\textbf{Memorylessness:} exponential waiting has the same residual distribution after any elapsed time. Uniform and general distributions do not; condition and renormalise their remaining support.
	\boxend

	\boxstart{L6 — Conditional survival templates}
	For a nonnegative duration $X$, translate verbal information into a survival event before calculating.\
	\textbf{Uniform:} conditioning on $X>s$ truncates the interval; the remaining time is uniform on the shortened interval.\
	\textbf{Exponential:} conditioning on $X>s$ leaves the residual time distributed as a fresh exponential variable.\
	\textbf{General distribution:} use
	\[
		P(X>s+t\mid X>s)=\frac{P(X>s+t)}{P(X>s)}.
	\]
	Do not invoke memorylessness for uniform, normal, Pareto, or arbitrary distributions.
	\boxend

	\boxstart{L6 — Waiting vs restarting}
	After surviving to time $s$, compare the expected remaining time $r(s)=E(X-s\mid X>s)$ with a fresh attempt mean $E(X)$. Restarting is better when $r(s)>E(X)$ (assuming no restart cost and an independent fresh attempt).

	\textbf{Pareto lecture example:} if $f_X(x)=2/x^3$ for $x\ge1$, then $P(X>s)=1/s^2$ and
	\[
		E(X\mid X>s)=\frac{\int_s^\infty x(2/x^3)\,dx}{1/s^2}=2s,
		\qquad r(s)=2s-s=s.
	\]
	Since $E(X)=2$, continue for $s\le2$ and restart for $s>2$.
	For a timeout $T>1$, the Pareto lecture example gives
	\[
		m(T)=\frac{2-1/T}{1-1/T^2},\qquad
		m'(T)=0\Longrightarrow T=\sqrt{2+\sqrt3}\approx3.73.
	\]
	\boxend

	\boxstart{L6 — Repeated timeout identity}
	Assume each unfinished attempt is abandoned at timeout $T$, then a fresh independent attempt starts immediately; no restart cost. If $m(T)$ is the expected total elapsed time,
	\[
		m(T)=P(X\le T)E(X\mid X\le T)+P(X>T)T+P(X>T)m(T),
	\]
	so
	\[
		\boxed{m(T)=\frac{E[\min(X,T)]}{P(X\le T)}}.
	\]
	For $X\sim\operatorname{Exp}(\lambda)$, $E[\min(X,T)]=(1-e^{-\lambda T})/\lambda$, hence $m(T)=1/\lambda$: repeated timeouts do not change the expected completion time.
	\boxend

	\boxstart{L5/L6 — Approximation decision rules}
	\textbf{Binomial $\to$ normal:} check both expected successes and failures are sufficiently large; use $Y\sim N(np,npq)$ and a continuity correction.
	\begin{itemize}
		\item $P(X=k)$ becomes $P(k-0.5<Y<k+0.5)$.
		\item $P(X\le c)$ becomes $P(Y<c+0.5)$; $P(X\ge c)$ becomes $P(Y>c-0.5)$.
	\end{itemize}
	\textbf{Binomial / Poisson-binomial $\to$ Poisson:} individual probabilities should be small; use mean $\lambda=\sum p_i$ (or $np$). Check whether the requested tail is sensitive to approximation error.
	\boxend

	\boxstart{L6 — Normal calculation template}
	For $X\sim N(\mu,\sigma^2)$, translate every bound using $Z=(X-\mu)/\sigma$ before using the supplied table:
	\[
		P(a<X<b)=\Phi\!\left(\frac{b-\mu}{\sigma}\right)-\Phi\!\left(\frac{a-\mu}{\sigma}\right).
	\]
	Use symmetry $\Phi(-z)=1-\Phi(z)$. For a central probability,
	\[
		P(|Z|<z)=2\Phi(z)-1.
	\]
	For a binomial count, apply continuity correction first, then standardise the corrected endpoints.
	\boxend

	\boxstart{L5/L6 — Poisson-process applications}
	For rate $r$ per unit time, count events in duration $t$ using parameter $rt$; counts in disjoint intervals are independent. Interarrival times are iid exponential with rate $r$.

	\textbf{Rate scaling:} per minute $\to$ multiply by minutes; per hour $\to$ multiply by hours. For multiple independent streams, add rates. For ``at least one'', use the zero-count complement.

	\textbf{Packet retransmission template:} if one attempt succeeds with probability $s$, attempts through one success are geometric with mean $1/s$. For $k$ accepted packets, total attempts is negative binomial; multiply expected attempts by the packet size to obtain expected transmitted bits.
	\boxend

	\boxstart{L6 — Simulation recipes from lecture}
	\textbf{Uniform point in unit disk:} draw independent $U,V\sim U(0,1)$; use $R=\sqrt U$, $\Theta=2\pi V$, then $(X,Y)=(R\cos\Theta,R\sin\Theta)$. Using $R=U$ is not uniform in area.

	\textbf{Independent Gaussian coordinates:} generate independent $U,V\sim U(0,1)$ and set
	\[
		R=\sqrt{-2\ln U},\quad\Theta=2\pi V,\quad
		Z_1=R\cos\Theta,\quad Z_2=R\sin\Theta.
	\]
	Then $Z_1,Z_2$ are independent $N(0,1)$; transform $X=\mu+\sigma Z$ for $N(\mu,\sigma^2)$.

	\textbf{Verification habit:} identify the target CDF/density, derive the distribution of the generated radius/transform, and check independence of inputs.
	\boxend

	\boxstart{L1–L6 — Exam traps and final checks}
	\begin{itemize}
		\item Do not use a distribution before checking its assumptions.
		\item Do not confuse a density with a probability or a CDF with a PMF.
		\item Do not omit a conditioning event from the denominator.
		\item Do not assume uncorrelated means independent.
		\item Do not forget support restrictions after marginalising or conditioning.
		\item For integer-valued tails, translate $>$, $\ge$, $<$, and $\le$ carefully before using a table.
		\item For normal approximations, standardise with the correct variance and apply the continuity correction once.
	\end{itemize}
	\textbf{Last pass:} units, support, probability range, edge cases, and whether all cases were counted exactly once.
	\boxend

\end{multicols*}
\end{document}
